\documentclass[10pt,a4paper]{amsart}
\usepackage[margin=19mm]{geometry}
\usepackage{amsmath,amssymb,mathtools}
\usepackage[scr=rsfs,cal=euler]{mathalfa}
\usepackage{xfrac}
\usepackage[unicode,hidelinks,bookmarks=false]{hyperref}
\newcommand{\ie}{i.e.\ }
\newcommand{\cf}{cf.\ }
\newcommand{\resp}{respectively\ }
\newcommand{\Subsection}{\subsection}
\providecommand{\nobreakdash}{\nobreak\hbox{-}}
\setlength{\emergencystretch}{2em}
\multlinegap0pt
\usepackage{amssymb}
\usepackage{dsfont}
\usepackage{mathtools}
\usepackage[shortlabels]{enumitem}
\newtheorem{proposition}{Proposition}
\DeclareMathOperator{\Pic}{Pic}
\newcommand{\Q}{\mathbb{Q}}
\newcommand{\Z}{\mathbb{Z}}
\newcommand{\defeq}{\vcentcolon=}
\title{Interpolation of generalized Heegner classes along quaternionic Coleman families}
\author{Eris Rocha Walchek}
\date{Source: arXiv:2503.01749v2 (CC BY 4.0)}
\begin{document}
\maketitle

\noindent\emph{Source and licence.} Interpolation of generalized Heegner classes along quaternionic Coleman families. Version arXiv:2503.01749v2 (CC BY 4.0), distributed by arXiv under the Creative Commons Attribution 4.0 International licence, http://creativecommons.org/licenses/by/4.0/, as recorded in the arXiv licence metadata for this exact version and shown on the abstract page https://arxiv.org/abs/2503.01749v2. Full licence notice: \texttt{licenses/arxiv-2503.01749-CC-BY-4.0.txt}.\par\smallskip

\noindent\textit{Editorial scope.} Counted unit: the article's proposition prop.HeegnerPoint (source0.tex lines 328--330). The article's complete original proof of it is delivered with it (source0.tex lines 331--340, one \texttt{proof} environment of 10 lines). No mathematics is written, repaired or re-derived: the statement and the proof are the article's own lines, in the article's own notation, and no retained line is substituted or re-flowed.  No further source-local context is needed: every cross-reference of every retained line resolves inside the two delivered spans.  Nothing was omitted from any retained span, so the package carries no omission record.  No retained line contains a citation, so the wrapper carries no bibliography and no bibliography entry.  No retained line contains a figure environment, an image-inclusion command or drawing code, so no image is delivered and the package declares no asset dependency.  Editorial formatting only, in addition: an AMS \texttt{amsart} preamble that restates the article's own declarations and macros used by the retained lines, namely the environments \texttt{proposition} on separate counters, together with the standard packages those lines need.  The retained environments print in this document as Proposition 1 in order of appearance, whereas the article prints the same items with the numbering of its own full document.  The complete native source of the article as distributed in the arXiv e-print for this exact version is bundled unchanged as \texttt{sources/174161/source0.tex}.
\begin{proposition}\label{prop.HeegnerPoint}
For any $[\mathfrak{a}]\in\Pic(\mathcal{O}_{cp^n})$, the point $x_{cp^n,m}([\mathfrak{a}])\in X_m(H_{cp^n})$ is a Heegner point of conductor $cp^n$.
\end{proposition}

\begin{proof}
We check for each prime $\ell$ if, writing $\xi\defeq \xi_{cp^n}^{(\ell)}$ and $\iota_\xi\defeq \xi^{-1}\iota_K^B\xi$, we have \[\iota_\xi(K_\ell)\cap (U_m\otimes\Z_\ell) = \iota_\xi(\mathcal{O}_{cp^n}\otimes \Z_\ell).\]
For $\ell\nmid N^+cp$, both orders $U_m\otimes\Z_\ell$ and $\mathcal{O}_{cp^n}\otimes \Z_\ell$ are maximal, thus the image of the latter is contained in the earlier. Next we consider $\ell\mid N^+cp$;

For a split prime $\ell=\mathfrak{l}\bar{\mathfrak{l}}\neq p$ (case $\ell=p$ is similar), where $\mathfrak{l}$ is the prime corresponding to the fixed embedding $\iota_\ell\colon\overline{\Q}\hookrightarrow\overline{\Q}_\ell$ from the beginning of the chapter, we have that $K_\ell\defeq K\otimes_{\Q}\Q_\ell$ splits as $K_\mathfrak{l}\oplus K_{\bar{\mathfrak{l}}}\cong\Q_\ell e_{\mathfrak{l}}\oplus\Q_\ell e_{\bar{\mathfrak{l}}}$, where \[ e_\mathfrak{l}=\frac{1\otimes\vartheta-\vartheta\otimes1}{(\vartheta-\bar{\vartheta})\otimes 1}\hspace{15pt}\text{and}\hspace{15pt}e_{\bar{\mathfrak{l}}}=\frac{\vartheta\otimes 1-1\otimes\bar{\vartheta}}{(\vartheta-\bar{\vartheta})\otimes1}.\] One then computes that for each $(a+b\vartheta)\otimes x\in K\otimes_\Q\Q_\ell$,
\[\iota_\xi((a+b\vartheta)\otimes x)=\begin{pmatrix}xa+xb\vartheta & \ell^{-s}xb\delta \\ 0 & xa+xb\bar{\vartheta}\end{pmatrix}=(xa+xb\vartheta)\iota_\xi(e_\mathfrak{l})+(xa+xb\bar{\vartheta})\iota_\xi(e_{\bar{\mathfrak{l}}}),\]
showing that $\iota_\xi(K_\ell)\cap (U_m\otimes\Z_\ell)$ consists of elements of the form \[(xa+xb\vartheta)e_\mathfrak{l}+(xa+xb\bar{\vartheta})e_{\bar{\mathfrak{l}}}\in\Z_\ell e_\mathfrak{l}\oplus\Z_\ell e_{\bar{\mathfrak{l}}}\] with $xb\in\ell^s\Z_\ell$, thus of elements in $\mathcal{O}_{cp^n}\otimes_\Z\Z_p$;

For a non-split prime $\ell$, one computes that for each $(a+b\vartheta)\otimes x\in K_\ell=K\otimes_\Q\Q_\ell$, \[\iota_\xi((a+b\vartheta)\otimes x)=\begin{pmatrix}xa & \ell^{-s}xb \\ xb\vartheta\bar{\vartheta}\ell^s & x(a+b(\vartheta+\bar{\vartheta}))\end{pmatrix},\] so elements $xa+xb\vartheta$ in $\iota_\xi(K_\ell)\cap (U_m\otimes\Z_\ell)$ are of the form $xa\in\Z_\ell$ and $xb\in\ell^s\Z_\ell$. This shows that, for all $s'\leq 2s$, $\iota_\xi$ is an optimal embedding into an Eichler order of level $\ell^{s'}$, showing both the cases inert ($s'=s$, when $\ell\mid c^{-}$ or $\ell=p$ inert) and ramified ($s'=2s$, when $\ell=p$ is ramified).\qedhere
\end{proof}

\end{document}
